Analisar como técnicas de aprimoramento de LLMs, como \textit{Zero-shot} \cite{zero-shot} e RAG, influenciam a geração de respostas voltadas para a pecuária leiteira, com o intuito de auxiliar produtores leiteiros na tomada de decisões, na otimização de processos, especialmente aqueles relacionados ao manejo, fundamentando as respostas em artigos científicos de uma base de dados externa especializada.

Sendo assim o presente objetivo desse projeto é:

\begin{enumerate}
\item Realizar uma revisão sistemática de estratégias para a adaptação de LLMs a contextos de aplicação específicos;
\item Coletar e pré-processar dados relacionados ao gado leiteiro para a geração de um conjunto de documentos especializados;
\item Avaliar estratégias de RAG em LLMs utilizando publicações científicas da área (por exemplo, \textit{Journal of Dairy Science});
\item Avaliar as respostas às perguntas por meio de técnicas e métricas para análise de coerência e desempenho do LLM.processos, baseadas em artigos científicos de uma base de dados externa especializada.
Sendo assim o presente objetivo desse projeto é:
1. Realizar uma revisão sistemática de estratégias para a adaptação de LLMs a contextos

\end{enumerate}